%%%% Начало оформления заголовка - оставить без изменений !!! %%%%
\input{my_folder/task_settings}	% настройки - начало 


{\centering%
	
	\Ministry\\
	\SPbPU\\
	\institute
\par}
}\intervalS% завершает input

				\noindent
				\begin{minipage}{\linewidth}
				\vspace{\mfloatsep} % интервал 	
				\begin{tabularx}{\linewidth}{Xl}
					&УТВЕРЖДАЮ      \\
					&\HeadTitle     \\			
					&\underline{\hspace*{0.1\textheight}} \Head     \\
					&<<\underline{\hspace*{0.05\textheight}}>> \underline{\hspace*{0.1\textheight}} \approveYear г.  \\  
				\end{tabularx}
				\vspace{\mfloatsep} % интервал 	
				\end{minipage}

\intervalS{\centering\bfseries%

				ЗАДАНИЕ\\
				на выполнение %с 2020 года 
				%по выполнению % до 2020 года
				выпускной квалификационной работы


\intervalS\normalfont%

				студенту \uline{\AuthorFullDat{}, гр.~\group}


\par}\intervalS%
%%%%
%%%% Конец оформления заголовка  %%%%
 	
	
	
\begin{enumerate}[1.]
	\item Тема работы: {\expandafter \ulined <<Разработка модели глубокого обучения для прогнозирования продуктивности сельскохозяйственных растений по генетическим маркерам>>.}
	%\item Тема работы (на английском языке): \uline{\thesisTitleEn.} % вероятно после 2021 года
	\item Срок сдачи студентом законченной работы: \uline{\thesisDeadline.} 
	\item Обзор литературы по теме ВКР:
		\begin{enumerate}[–]
			\item Juan,W.;Ahn,K.W.;Chen,Y.-G.;Lin, C.-W. DropDAE:Denosing Autoencoder with Contrastive Learning for Addressing Dropout Events inscRNA-seqData. Bioengineering 2025, 12, 829. https://doi.org/10.3390/ bioengineering12080829
			
			\item Lu Q, Zhang M, Niu X, Wang S, Xu Q, Feng Y, Wang C, Deng H, Yuan X, Yu H, Wang Y, 
			Wei X. Genetic variation and association mapping for 12 agronomic traits in indica rice. 
			BMC Genomics. 2015 Dec 16;16:1067. doi: 10.1186/s12864-015-2245-2. PMID: 
			26673149; PMCID: PMC4681178.
			
			\item Jiaxin Chen, Cong Tan, Min Zhu, Chenyang Zhang, Zhihan Wang, Xuemei Ni, Yanlin Liu, 
			Tong Wei, XiaoFeng Wei, Xiaodong Fang, Yang Xu, Xuehui Huang, Jie Qiu, Huan Liu, 
			CropGS-Hub: a comprehensive database of genotype and phenotype resources for genomic 
			prediction in major crops, Nucleic Acids Research, Volume 52, Issue D1, 5 January 2024, 
			Pages D1519–D1529, https://doi.org/10.1093/nar/gkad1062
		\end{enumerate}
	\item Исходные данные по работе: 
	
	\uline{Данные генетических маркеров и фенотипических характеристик сельскохозяйственных 
		растений, полученные из открытых научных источников и специализированных баз данных..}% не забываем точку!
	
	\item Инструментальные средства:
	
	\begin{enumerate}[–]
		\item языки программирования: Python
		\item среда разработки: Visual Studio Code
		\item система контроля версий: git
	\end{enumerate}

	\printbibliographyTask % печать списка источников % КОММЕНТИРУЕМ ЕСЛИ НЕ ИСПОЛЬЗУЕТСЯ
	% В СЛУЧАЕ, ЕСЛИ НЕ ИСПОЛЬЗУЕТСЯ МОЖНО ТАКЖЕ ЗАЙТИ В setup.tex и закомментировать \vspace{-0.28\curtextsize}
	\item Содержание работы:
	\begin{enumerate}[label=\theenumi\arabic*.]
		\item Введение. Обоснование актуальности проблемы.
		\item Постановка задачи.
		\item Архитектура модели.
		\item Конфигурации модели и их применение. Вычислительные эксперименты.
		\item Результаты исследований.
		\item Выводы.
		\item Заключение. 
	\end{enumerate}
	\item Перечень графического материала: 
	\begin{enumerate}[label=\theenumi\arabic*.]
		\item Схема работы метода/алгоритма.%не забываем точку
		\item Архитектура разработанной программы/библиотеки.%не забываем точку
	\end{enumerate}	

	\item Консультанты по работе:
	\begin{enumerate}[label=\theenumi\arabic*.] 
	\item \uline{\ConsultantNorm{} (нормоконтроль).} %	Обязателен для всех студентов
	\end{enumerate}
		\item Дата выдачи задания: \uline{02.10.2025.}
\end{enumerate}

\intervalS%можно удалить пробел

Руководитель ВКР \uline{\hspace*{0.1\textheight} \Supervisor}


\intervalS%можно удалить пробел

Консультант  \uline{\hspace*{0.1\textheight}\ConsultantExtra}


\intervalS%можно удалить пробел

%Консультант по нормоконтролю \uline{\hspace*{0.1\textheight} \ConsultantNorm}%ПОКА НЕ ТРЕБУЕТСЯ, Т.К. ОН У ВСЕХ ПО УМОЛЧАНИЮ

Задание принял к исполнению \uline{\thesisStartDate}

\intervalS%можно удалить пробел

Студент \uline{\hspace*{0.1\textheight}  \Author}



\input{my_folder/task_settings_restore}	% настройки - конец